\begin{document}

{
\setcounter{tocdepth}{2}
\tableofcontents
}
\section{TFPT / Universalraum: gemeinsame Schatten und ein
ursprünglicher kausaler
Prozess}\label{tfpt-universalraum-gemeinsame-schatten-und-ein-urspruxfcnglicher-kausaler-prozess}

Vollständige Fortschreibung v1.6.8, 15. September 2026.
Forschungsbericht und unabhängige Prüfung der nachgereichten
Forschungssynthese 1.0. Keine vollständige TOE.

\subsection{1. Ergebnis und
Leseschlüssel}\label{ergebnis-und-leseschluxfcssel}

Die drei parallel bearbeiteten Aufgaben liefern diesmal konkrete
mathematische Anschlüsse: einen kleinen ursprünglichen
Singulettmechanismus, die vollständige bedingte Schattenkarte des hellen
Raums und einen kausalen Eingriff unter dem unveränderten
Hamiltonoperator. Die nachgereichte Forschungssynthese verbindet diese
Aufgaben sinnvoll durch gemeinsame Operationsgeschichten. Ihre
mitgelieferten Rechnungen sind reproduziert; ein Satz in Abschnitt 8.4
benötigt jedoch eine Korrektur.

Die noch später eingegangenen Texte liefern außerdem eine frisch
reproduzierte fünfte Bosonstufennorm. Eine unabhängig zertifizierte
größere Variationskompression verbessert die Grundenergieobergrenze auf
-1.13847609Δ und unter dem beibehaltenen N=63-Satz die untere
Polenergiegrenze auf 0.01657709Δ. Dazu kommt ein neuer exakter
Händigkeitstest des Feldprojektors. Der Nachtrag „Neue Norm, korrigierte
Zustandszuordnung und Feldadapter`` enthält diese maßgeblichen
Ergänzungen; ältere Ritzkennzahlen verschiedener Zustände werden dort
ausdrücklich getrennt.

Die zentrale Vereinfachung ist nicht eine neue große Geometrie, sondern
eine sorgfältig ausgewählte kleine gemeinsame Ausführung. Ein Zustand,
seine messbaren Bilder und seine späteren Antworten müssen aus denselben
Operatoren berechnet werden. Zusätzliche Zustände, Kontrollen und
Referenzen werden dabei ausdrücklich mitgezählt.

\textbf{Exakt} bezeichnet algebraische Identitäten, Invarianzbeweise und
rational/modular abgesicherte Ränge im angegebenen Modell.
\textbf{Numerisch} bezeichnet unabhängige Gleitkomma-Rekonstruktionen
und den vollständigen endlichen Dynamikvergleich. \textbf{Bedingt}
heißt: Der Satz stimmt, benötigt aber ausdrücklich gewährte
Präparationen, Kontrollen oder Messungen. \textbf{Offen} bleibt deren
ursprüngliche Herleitung und die physikalische Fortsetzung.
Prüfprogramme sind keine formale Beweisassistenten-Verifikation.

Die vollständige bisherige Hauptfassung v1.6.7 bleibt im historischen
Teil wortgetreu erhalten. Die neuen Kapitel, die vollständigen drei
Teilberichte, ihr gemeinsames Dreizustands-Addendum und alle drei
Nachtragsprüfungen stehen davor. Das Quellenpaper bleibt als
unveränderte Eingabe im Prüfpaket, nicht als still korrigierte
Originaldatei. Aktuell maßgeblich sind diese Einordnung, die Korrekturen
und die neu dokumentierten Ressourcenbedingungen.

\subsection{2. Prüfung des nachgereichten
Forschungspapers}\label{pruxfcfung-des-nachgereichten-forschungspapers}

Geprüft wurden das 33-seitige PDF, sein vollständiger Markdowntext,
Prüfprotokoll, Quellenmanifest und eigene Ergebnisdatei sowie das
zusätzlich vorhandene Reproduktionsarchiv. Die vollständige Synthese
wurde gelesen; nicht jede der 74 historischen Quellen wurde neu
bewiesen. Die relevanten PDF-Seiten 16 bis 20 und 22 wurden zusätzlich
bildlich geprüft.

{\def\LTcaptype{none} % do not increment counter
\begin{longtable}[]{@{}
  >{\raggedright\arraybackslash}p{(\linewidth - 2\tabcolsep) * \real{0.5000}}
  >{\raggedright\arraybackslash}p{(\linewidth - 2\tabcolsep) * \real{0.5000}}@{}}
\toprule\noalign{}
\begin{minipage}[b]{\linewidth}\raggedright
Behauptung oder Baustein
\end{minipage} & \begin{minipage}[b]{\linewidth}\raggedright
Ergebnis der unabhängigen Prüfung
\end{minipage} \\
\midrule\noalign{}
\endhead
\bottomrule\noalign{}
\endlastfoot
74 eingefrorene Quellen & Alle Archivgrößen und SHA-256-Werte stimmen
mit dem Manifest überein \\
Markdown, PDF und TeX & Gelieferte bzw. archivierte Dateien stimmen mit
dem Prüfprotokoll überein \\
Neuester Stand S074 & Identisch mit der vollständigen hiesigen v1.6.7
einschließlich Prozess-/Schattennachtrag \\
57 eigene Prüfbedingungen & Aus dem isoliert entpackten Paket normal und
optimiert wiederholt; identische Ergebnisbytes \\
20 übernommene Kettenprüfungen & Ebenso wiederholt; identische
Ergebnisbytes \\
Grundzustand und frühere 822er-Suite & In dieser Runde nicht vollständig
erneut ausgeführt; keine entsprechende Behauptung \\
Wortkernrekonstruktion & Unter den angegebenen Positivitäts-, Relations-
und Beschränktheitsvoraussetzungen korrekt \\
Aussage zu der ersten Gramzeile in §8.4 & Zu korrigieren; siehe unten \\
Ein-Cartan-Referenz und Projektorgeometrie & Korrekte bedingte
Konstruktionen, kein nativer physikalischer Ursprung \\
\end{longtable}
}

\subsubsection{2.1 Tatsächliche Korrektur: die erste Zeile genügt, wenn
sie vollständig
ist}\label{tatsuxe4chliche-korrektur-die-erste-zeile-genuxfcgt-wenn-sie-vollstuxe4ndig-ist}

Für eine unter Produkt und Adjungieren abgeschlossene Wortalgebra gilt
definitionsgemäß

\[\adjustbox{max width=\linewidth}{$\displaystyle 
\boxed{\Gamma(u,v)=\omega(u^\dagger v)=\Gamma(1,u^\dagger v).}
$}\]

Damit bestimmt die Gesamtheit aller Werte \(\Gamma(1,w)\) bereits den
gesamten Gramkern. Der Satz in §8.4, Gleichheit nur von \(\Gamma(1,u)\)
genüge nicht, ist ohne Einschränkung auf unvollständige Daten falsch.
Die richtige Fassung lautet:

\begin{quote}
Gleichheit nur endlich vieler Worterwartungswerte, einzelner
Zeitkorrelationen oder Ausgangswahrscheinlichkeiten genügt im
Allgemeinen nicht zur Identifikation des vollständigen Prozesses.
\end{quote}

Die Korrektur betrifft nicht den GNS-Rekonstruktionssatz. Sie zeigt
vielmehr: Die zweifach indizierte Grammatrix ist eine hilfreiche
Organisation derselben vollständigen Daten, keine zusätzlich notwendige
unabhängige Informationsquelle. Die Identität wird im eigenen
Ressourcenprüfer auf einem nichtkommutativen Matrixbeispiel zusätzlich
kontrolliert; der allgemeine Beweis ist die eine Gleichung oben. Die
gelieferten 57 Prüfbedingungen hatten die fehlerhafte Formulierung nicht
geprüft.

\subsubsection{2.2 Welche Gleichheit ist tatsächlich
beobachtbar?}\label{welche-gleichheit-ist-tatsuxe4chlich-beobachtbar}

Ein markierter komplexer Wortkern ist stärker als eine Liste
gewöhnlicher Quantenkanäle. Beispielsweise liefern \(Z\) und \(-Z\) für
jeden Zustand denselben Kanal:

\[\adjustbox{max width=\linewidth}{$\displaystyle 
Z\rho Z^\dagger=(-Z)\rho(-Z)^\dagger.
$}\]

Mit \(\Omega=|0\rangle\) sind die markierten Wortwerte dagegen +1 und
-1. Bei verfügbarer kohärenter Kontrolle der Zweige können diese Phasen
relativ werden und verschiedene Interferenz ergeben. Ohne diese
zusätzliche Kontrollmöglichkeit ist ihre Unterscheidung kein
zugängliches Experiment. Diese Ressourcenfrage ist auch aus der
Literatur zur Kontrolle unbekannter Operationen bekannt; unser konkreter
Zwei-Zustandszeuge wird unabhängig direkt nachgerechnet.
\href{https://arxiv.org/abs/1309.7976}{Araújo, Feix, Costa und Brukner,
2014}.

Folglich ist Kernelgleichheit ein hinreichender starker
Identifikationsvertrag für phasenmarkierte gemeinsame Ausführungen. Sie
ist nicht ohne Weiteres die notwendige Minimalbeschreibung jedes
eingeschränkten Instrumentenvertrags. Man darf weder beobachtbare
Interferenz wegquotientieren noch unzugängliche globale Kanalphasen als
neue physikalische Unterschiede ausgeben.

\subsubsection{2.3 Referenzsystem: richtige Bilanz, keine kostenlose
Operation}\label{referenzsystem-richtige-bilanz-keine-kostenlose-operation}

Die Synthese ergänzt den früher zusätzlich gemischten Vierzustandskanal
um eine zweistufige Referenz. Das acht-dimensionale Gesamtmodell erhält
die betrachtete innere Cartanladung. Auf einem invarianten
vierdimensionalen Unterraum ist sein Hamiltonoperator exakt der alte
Transferblock plus ein skalares Energieglied. Die frühere
Transferschranke wird deshalb korrekt übertragen, nicht neu aus der
nativen Quelle erzeugt.

Der gewählte innere Cartanwert ist nicht die physische Zahl
\(N=N_f+2N_b\). Die gesamte nichtabelsche Gruppe bleibt in dieser
Konstruktion ungeprüft. Der Referenzwechsel wird bezahlt: Scharfe
verschiedene Referenzladungen sind orthogonal. Nach Wegspuren der
Referenz fehlen die Kreuzterme zwischen verschiedenen Systemladungen.
Die gemeinsame kohärente Entwicklung ist somit kein unverändert
wiederverwendbarer kohärenter Antrieb des isolierten Systems. Das
Original benennt diese Grenzen bereits; sie werden bei der Integration
beibehalten.

\subsubsection{2.4 Projektorgeometrie: tragfähig, aber noch nicht
nativ}\label{projektorgeometrie-tragfuxe4hig-aber-noch-nicht-nativ}

Für eine Isometriefamilie J mit \(P=JJ^\dagger\) gilt

\[\adjustbox{max width=\linewidth}{$\displaystyle 
A=iJ^\dagger dJ,\qquad
F=dA-iA\wedge A=i\,dJ^\dagger(I-P)\wedge dJ.
$}\]

Das verbindet Quantenmetrik und geometrische Phase eines wirklich
bewegten Unterraums. Eine reine Basisdrehung im festen vollständigen
Raum erzeugt keine solche Krümmung. Der im Paper gerechnete Kugelzeuge
mit \(u=(\cos(\theta/2),e^{i\phi}\sin(\theta/2))\) und seiner selektiven
Schleifenausbeute 1/8 stimmt. Die Parameterfamilie wird jedoch
hinzugegeben. Die Projektorformulierung entspricht etablierten Verfahren
der Quantengeometrie. \href{https://arxiv.org/abs/2102.09899}{Graf und
Piéchon, 2021}.

Eine neue genaue Einschränkung folgt aus unserem ursprünglichen hellen
Block: Sein Blochvektor hat bei reellen Kopplungen nur zwei Komponenten,
\(n\propto(2a,0,-\Delta)\). Daher ist

\[\adjustbox{max width=\linewidth}{$\displaystyle 
n\cdot(\partial_a n\times\partial_\Delta n)=0.
$}\]

Die glatte lokale Berry-Zweiform dieser realen Zweiparameterfamilie
verschwindet außerhalb einer Entartung. Eine komplexe Kontrollphase oder
eine andere tatsächlich bewegte Projektorfamilie wäre eine zusätzliche,
noch herzuleitende Ressource. Globale Vorzeichenphasen auf anderen
Parametergebieten werden dadurch nicht allgemein ausgeschlossen; eine
gekrümmte Raumzeit folgt weder aus dem Kugelbeispiel noch aus diesem
Nullbefund.

\subsubsection{2.5 Weitere übernommene
Richtungen}\label{weitere-uxfcbernommene-richtungen}

Die Blockresolventenformel und der endliche Duhamel-Adapterfehlervertrag
sind korrekt und brauchbar. Die einfache Restschranke
\(\|B\|^2/|\operatorname{Im}z|^3\) wird nahe der reellen Achse schwach;
sie bestimmt keinen scharfen physikalischen Pol. Gemischte Wort- und
Matrixmomente sind eine passende Fortsetzung der verzweigten
Hamiltonkette. Das neue Paper hat diese größere Kontraktionsrechnung
selbst noch nicht ausgeführt. Es liefert auch kein neues chirales Maß,
keinen Spin-2-Sektor und keinen gemeinsamen 3+1D-Ursprung. Seine
Zurückhaltung an diesen Stellen ist sachlich richtig.

\subsection{3. Der gemeinsame ursprüngliche
Ausgangspunkt}\label{der-gemeinsame-urspruxfcngliche-ausgangspunkt}

Alle drei eigenen Stränge verwenden denselben eingefrorenen Tensor W mit
60 Zeilen, 2016 Paarspalten und 480 signierten Einträgen. Es gilt exakt

\[\adjustbox{max width=\linewidth}{$\displaystyle 
WW^\dagger=8I_{60},\quad
P_A=\sum_{i<j}W_{A,ij}f_jf_i,\quad
H=\Delta N_b+g\sum_A(b_A^\dagger P_A+P_A^\dagger b_A),
\quad N=N_f+2N_b.
$}\]

Der Dateipin ist
\codeword{3f00a0892157a691e4ef26920c3702772c7bf76a250fad98ff04a079bb64b763}.
Der ursprüngliche Clifford-/Farbkonstruktor wurde im Mechanismusstrang
erneut ausgeführt und reproduziert W eintragsweise. Die endliche
Clockkonstruktion wurde in ihren vier angegebenen mathematischen Stufen
wiederholt. Diese Wiederholung liefert die konkreten Koeffizienten und
den passiven Lift; sie wählt weder die physikalische Kopplung noch
Apparate oder eine Raumzeit.

Die neue Arbeit betrifft mehrere sauber bezeichnete Teilräume. Gleiche
Quelle heißt nicht gleicher Zustand:

{\def\LTcaptype{none} % do not increment counter
\begin{longtable}[]{@{}
  >{\raggedright\arraybackslash}p{(\linewidth - 4\tabcolsep) * \real{0.3333}}
  >{\raggedright\arraybackslash}p{(\linewidth - 4\tabcolsep) * \real{0.3333}}
  >{\raggedright\arraybackslash}p{(\linewidth - 4\tabcolsep) * \real{0.3333}}@{}}
\toprule\noalign{}
\begin{minipage}[b]{\linewidth}\raggedright
Teilraum
\end{minipage} & \begin{minipage}[b]{\linewidth}\raggedright
Inhalt
\end{minipage} & \begin{minipage}[b]{\linewidth}\raggedright
Wofür benutzt
\end{minipage} \\
\midrule\noalign{}
\endhead
\bottomrule\noalign{}
\endlastfoot
Heller Raum bei N=2, Dimension 120 & Kohärente Fermionpaare und ein
Boson, je 60 innere Richtungen & Gemeinsame Schatten- und
Zeitrekonstruktion \\
Ein nativer Stern bei N=2, Dimension 9 & Acht disjunkte Paare und ein
gemeinsames Boson & Kausaler Phasenimpuls, einschließlich dunkler
Richtungen \\
Neuer Singulett-Unterraum bei N=4, Dimension 2 & Zwei Bosonen bzw. ein
Boson und zwei Fermionen & Kleinster hier explizit geschlossener voller
G-Singulettmechanismus \\
N=4 mit neuer Referenz, Dimension 3 & Vakuum/Referenz plus die beiden
Singulettzustände & Bedingter neutraler Ressourcenübergang \\
Früherer Grundzustand bei N=64 & Native geladene Polantwort & Nicht
durch die neuen N=2-/N=4-Zeugen ersetzt \\
\end{longtable}
}

Der Stern enthält sieben dunkle Paarrichtungen und ist nicht vollständig
im hellen 120-Raum enthalten. Ein lokaler Phasenimpuls kann helle in
dunkle Komponenten überführen. Deshalb darf die helle Tomographie nicht
ohne weiteren Adapter als vollständige Tomographie des kausalen
Sternversuchs ausgegeben werden.

\subsection{4. Einfachster Mechanismus: zwei Zustände und eine genaue
Operationsgrenze}\label{einfachster-mechanismus-zwei-zustuxe4nde-und-eine-genaue-operationsgrenze}

Mit der aus W hergeleiteten invarianten Paarung \(\eta\) definiere

\[\adjustbox{max width=\linewidth}{$\displaystyle 
B_+=\tfrac12 b^\dagger\eta b^\dagger,\qquad
R_+=\sum_{AB}\eta_{AB}b_A^\dagger P_B^\dagger.
$}\]

Die vollständige Anwendung der nativen Terme ergibt

\[\adjustbox{max width=\linewidth}{$\displaystyle 
\sum_{AB}\eta_{AB}P_A^\dagger P_B^\dagger=0,
\quad T_-B_+|0\rangle=R_+|0\rangle,
\quad T_+R_+|0\rangle=16B_+|0\rangle.
$}\]

Die jeweils anderen Übergänge sind null. Die erste Identität wurde durch
Normalordnung sämtlicher 3840 signierten Beiträge auf 960
Fermionquartetten geprüft. Mit

\[\adjustbox{max width=\linewidth}{$\displaystyle 
\beta=B_+|0\rangle/\sqrt{30},\quad
\rho=R_+|0\rangle/\sqrt{480}
$}\]

ist der ganze H auf diesem invarianten Unterraum genau

\[\adjustbox{max width=\linewidth}{$\displaystyle 
\boxed{H_4=\begin{pmatrix}2\Delta&4g\\4g&\Delta\end{pmatrix}.}
$}\]

Beide Zustände haben N=4 und sind volle innere G-Singuletts. Keine
höhere Konfiguration wurde weggelassen. Am Prüfpunkt \(g/\Delta=1/20\)
beträgt die maximale Umwandlungswahrscheinlichkeit exakt 4/29. Ihre
Vorbereitung aus dem Vakuum ist dadurch nicht bewiesen.

Tatsächlich folgt für das dokumentierte Alphabet aus passiven Focklifts,
X und Nb ein geschlossener Erhaltungssatz:

\[\adjustbox{max width=\linewidth}{$\displaystyle 
\mathcal A_0\subseteq\{N\}'.
$}\]

Er gilt für beliebige Produkte, Adjungierte und Kommutatoren; weitere
Wortlängen können ihn nicht umgehen. B+ und R+ ändern N um vier und
liegen nicht in dieser Algebra. Der Ausschluss gilt für dieses geprüfte
Alphabet, nicht für jede denkbare alternative Compilerrealisierung.

Eine minimale zusätzliche G-invariante, fermionparitätsgerade Operation
ist nach Polynomialgrad \(Q=B_++B_-\). Unter zusätzlich separat
zugänglichen X und Nb gilt exakt

\[\adjustbox{max width=\linewidth}{$\displaystyle 
\boxed{R_++R_-=-[N_b,[Q,X]].}
$}\]

Man braucht also nicht zwei neue unabhängige Kopplungstensoren. Ein
zusätzlicher Paargriff reicht algebraisch. Das ist keine bereits
ausgeführte endliche Pulsynthese und keine Herkunft von Q.

Eine zweistufige, innerlich G-triviale Referenz mit N-Ladungen 0 und 4
bilanziert diesen einen zusätzlichen Griff. Auf
\(|0,4_R\rangle,|\beta,0_R\rangle,|\rho,0_R\rangle\) folgt der exakte
gemeinsame Block

\[\adjustbox{max width=\linewidth}{$\displaystyle 
H_{\rm ref}=\begin{pmatrix}
E_R&\sqrt{30}\kappa&0\\
\sqrt{30}\kappa&2\Delta&4g\\
0&4g&\Delta
\end{pmatrix}.
$}\]

Er erhält volle innere G-Symmetrie und Gesamt-N=4. Referenz, Energie ER
und Kopplung κ sind neue Ressourcen. Er löst ein anderes Ladungsproblem
als die Ein-Cartan-Referenz im gelieferten Paper und erzeugt noch keinen
räumlichen Transfer. Sein vollständiger Beweis und die reduzierte
Zustandsbilanz stehen im Mechanismus-Teilbericht.

